\subsection*{给频率测量增加一个衰减旋钮}
傅里叶测量把信号乘上反向旋转，再累积。如果信号在无穷远持续增长，积分可能不收敛。拉普拉斯核加入指数衰减：写 $s=\sigma+\mathrm i\omega$，核为 $\mathrm e^{-\sigma t}\mathrm e^{-\mathrm i\omega t}$。一个旋钮控制遗忘远处的速度，另一个控制匹配频率。

对 $f(t)=\mathrm e^{at}$，乘核后是 $\mathrm e^{-(s-a)t}$。只有 $\operatorname{Re}s>a$，衰减才压过增长，积分得到 $1/(s-a)$。收敛条件来自积分的尾部，不是公式之外可有可无的小注释。

\subsection*{一张有理式不能代替全部变换信息}
令 $a=2$，在实数 $s=3$ 处，原积分收敛且为一；在 $s=1$，有理式形式上给负一，但原积分发散。这个例子说明解析表达式延伸到某处有值，不意味着它仍代表定义积分。

单边变换从时间零开始，适合记录带初值的过程。选择单边或双边、选择何种收敛域，会影响反演与解释；因此每次使用变换表，都应同时读变换约定和适用区域。

这一节先把变换理解成一族带权测量。下一节真正有用的地方是：求导在这种测量下转为乘以 $s$，并留下起点信息。原本关于时间函数的微分方程，因此可能变成关于变换值的代数关系。
